\section{Tranche 4: Hierarchical Network Construction}
\label{sec:network_model}

In the final tranche, the validated JSON trees from Tranche~3 are instantiated into a longitudinal \textbf{Hierarchical Network of Networks} across the five annual snapshots $t \in \{2021, 2022, 2023, 2024, 2025\}$. For each year $t$, the multi-scale system is formally defined as:
\begin{equation}
    \Hnet = \left( \Gmacro, \; \left\{ \Gmicro \right\}_{i=1}^{N} \right), \quad N = 108
\end{equation}
where $\Gmacro$ denotes the macroscopic inter-firm network over the static cohort of $N=108$ QQQ firms, and each macro-node $C_i \in \Vmacro$ encapsulates a rooted intra-firm micro-network $\Gmicro$ (Fig.~\ref{fig:micro_macro_schema}).

\begin{figure}[p]
\centering
\resizebox{0.95\textwidth}{!}{%
\begin{tikzpicture}[
    >=Stealth,
    % High-contrast, 3-shade B&W print-safe palette
    rootnode/.style={ellipse, draw=black!95, double, double distance=1.3pt, thick, fill=gray!25, minimum width=2.4cm, minimum height=0.95cm, align=center, font=\small\bfseries},
    hubnode/.style={ellipse, draw=black!90, very thick, fill=gray!14, minimum width=2.05cm, minimum height=0.82cm, align=center, font=\footnotesize\bfseries},
    subnode/.style={ellipse, draw=black!80, semithick, fill=white, minimum width=1.55cm, minimum height=0.62cm, align=center, font=\scriptsize},
    segnode/.style={ellipse, draw=black!90, thick, fill=white, minimum width=1.6cm, minimum height=0.62cm, align=center, font=\scriptsize\bfseries},
    regnode/.style={rectangle, rounded corners=6pt, draw=black!85, thick, fill=white, minimum width=1.45cm, minimum height=0.58cm, align=center, font=\scriptsize\bfseries},
    prodleaf/.style={ellipse, draw=black!80, semithick, fill=gray!8, minimum width=2.15cm, minimum height=0.76cm, align=center, font=\scriptsize},
    dotleaf/.style={circle, draw=black!90, semithick, fill=black!50, inner sep=2.3pt},
    macronode/.style={circle, draw=black!85, thick, fill=white, align=center, font=\scriptsize\bfseries},
    annot/.style={font=\footnotesize\itshape, align=left, text=black!90},
    edgestyle/.style={draw=black!85, semithick},
    dotedge/.style={draw=black!85, very thick, dashdotted}
]

% =========================================================================
% 1. MICRO-LEVEL TOPOLOGY (TOP: COMPANY 8)
% =========================================================================

% --- Center Root Node ---
\node[rootnode] (comp8) at (0, 0) {Company 8 [1]};

% -------------------------------------------------------------------------
% A. TOP BRANCH: MARKET COMPETITION [4, 5]
% -------------------------------------------------------------------------
\node[hubnode] (comp_hub) at (0, 1.9) {Market\\Competition};
\node[subnode] (comp_seg) at (-1.6, 3.2) {Segment [4]};
\node[subnode] (comp_oth) at (1.8, 3.1) {Others};

\node[segnode] (comp_sa) at (-3.5, 4.5) {$\text{Segment}_a$};
\node[segnode] (comp_sb) at (-1.6, 4.7) {$\text{Segment}_b$};
\node[segnode] (comp_sn) at (0.3, 4.5) {$\text{Segment}_n$};

\draw[edgestyle] (comp8) -- (comp_hub);
\draw[edgestyle] (comp_hub) -- (comp_seg);
\draw[edgestyle] (comp_hub) -- (comp_oth);
\draw[edgestyle] (comp_seg) -- (comp_sa);
\draw[edgestyle] (comp_seg) -- (comp_sb);
\draw[edgestyle] (comp_seg) -- (comp_sn);

% Competitor leaf nodes
\node[dotleaf] (csa1) at (-4.2, 5.4) {};
\node[dotleaf] (csa2) at (-3.5, 5.6) {};
\node[dotleaf] (csa3) at (-2.8, 5.4) {};
\draw[edgestyle] (comp_sa) -- (csa1);
\draw[edgestyle] (comp_sa) -- (csa2);
\draw[edgestyle] (comp_sa) -- (csa3);

\node[dotleaf] (csb1) at (-2.0, 5.6) {};
\node[dotleaf] (csb2) at (-1.2, 5.6) {};
\draw[edgestyle] (comp_sb) -- (csb1);
\draw[edgestyle] (comp_sb) -- (csb2);

\node[dotleaf] (csn1) at (-0.1, 5.4) {};
\node[dotleaf] (csn2) at (0.6, 5.5) {};
\node[dotleaf] (csn3) at (1.1, 5.1) {};
\draw[edgestyle] (comp_sn) -- (csn1);
\draw[edgestyle] (comp_sn) -- (csn2);
\draw[edgestyle] (comp_sn) -- (csn3);

\node[dotleaf] (coth1) at (2.7, 3.9) {};
\node[dotleaf] (coth2) at (3.0, 3.1) {};
\node[dotleaf] (coth3) at (2.6, 2.3) {};
\draw[edgestyle] (comp_oth) -- (coth1);
\draw[edgestyle] (comp_oth) -- (coth2);
\draw[edgestyle] (comp_oth) -- (coth3);

\node[annot, align=center] at (-1.5, 6.2) {\textbf{Competitor Nodes [5]}\\\textit{Explicit \& Inferred}};

% -------------------------------------------------------------------------
% B. RIGHT BRANCH: DUAL-SCOPE RISK TAXONOMY [6, 7, 8]
% -------------------------------------------------------------------------
\node[hubnode] (risk_hub) at (3.5, 0.1) {Dual-Scope\\Risk Taxonomy};
\node[subnode] (risk_seg) at (5.0, 1.5) {Categories [6]};
\node[subnode] (risk_oth) at (5.1, -1.3) {Unified Risks [7]};

\node[segnode] (risk_sa) at (7.2, 2.7) {$\text{Risk}_a$};
\node[segnode] (risk_sb) at (7.4, 1.5) {$\text{Risk}_b$};
\node[segnode] (risk_sn) at (7.2, 0.3) {$\text{Risk}_n$};

\draw[edgestyle] (comp8) -- (risk_hub);
\draw[edgestyle] (risk_hub) -- (risk_seg);
\draw[edgestyle] (risk_hub) -- (risk_oth);
\draw[edgestyle] (risk_seg) -- (risk_sa);
\draw[edgestyle] (risk_seg) -- (risk_sb);
\draw[edgestyle] (risk_seg) -- (risk_sn);

% Risk leaf nodes
\node[dotleaf] (rsa1) at (8.5, 3.2) {};
\node[dotleaf] (rsa2) at (8.7, 2.7) {};
\node[dotleaf] (rsa3) at (8.4, 2.2) {};
\draw[edgestyle] (risk_sa) -- (rsa1);
\draw[edgestyle] (risk_sa) -- (rsa2);
\draw[edgestyle] (risk_sa) -- (rsa3);

\node[dotleaf] (rsb1) at (8.8, 1.8) {};
\node[dotleaf] (rsb2) at (8.8, 1.2) {};
\draw[edgestyle] (risk_sb) -- (rsb1);
\draw[edgestyle] (risk_sb) -- (rsb2);

\node[dotleaf] (rsn1) at (8.5, 0.7) {};
\node[dotleaf] (rsn2) at (8.7, 0.2) {};
\node[dotleaf] (rsn3) at (8.3, -0.3) {};
\draw[edgestyle] (risk_sn) -- (rsn1);
\draw[edgestyle] (risk_sn) -- (rsn2);
\draw[edgestyle] (risk_sn) -- (rsn3);

\node[annot, align=center] at (8.8, 4.0) {\textbf{Risk Drivers}\\\textit{Operational Triggers}};

\node[dotleaf] (roth1) at (6.6, -1.0) {};
\node[dotleaf] (roth2) at (6.4, -1.9) {};
\draw[edgestyle] (risk_oth) -- (roth1);
\draw[edgestyle] (risk_oth) -- (roth2);

% MOVED M:N Bridge text far right into open space to prevent collision
\node[annot, align=left] at (8.5, -1.6) {\textbf{M:N Associative Bridge [8]}\\(Multi-segment threats)};

% -------------------------------------------------------------------------
% C. LEFT BRANCH: WRDS GEOGRAPHIC FINANCIALS [10]
% -------------------------------------------------------------------------
\node[hubnode] (geo_hub) at (-3.5, 0.4) {WRDS Geographic\\Financials [10]};
\node[subnode] (geo_sub) at (-5.4, 1.8) {Geographic};
\node[subnode] (geo_oth) at (-5.6, -0.7) {Other};

\node[regnode] (reg_a) at (-7.9, 2.3) {$\text{Region}_a$};
\node[regnode] (reg_b) at (-7.1, 3.4) {$\text{Region}_b$};
\node[regnode] (reg_n) at (-5.4, 3.8) {$\text{Region}_n$};
\node[regnode, minimum width=1.2cm] (reg_rest) at (-3.9, 2.8) {Rest};

\draw[edgestyle] (comp8) -- (geo_hub);
\draw[edgestyle] (geo_hub) -- (geo_sub);
\draw[edgestyle] (geo_hub) -- (geo_oth);
\draw[edgestyle] (geo_sub) -- (reg_a);
\draw[edgestyle] (geo_sub) -- (reg_b);
\draw[edgestyle] (geo_sub) -- (reg_n);
\draw[edgestyle] (geo_sub) -- (reg_rest);

\node[dotleaf] (goth1) at (-6.7, -1.6) {};
\node[dotleaf] (goth2) at (-5.5, -1.8) {};
\draw[edgestyle] (geo_oth) -- (goth1);
\draw[edgestyle] (geo_oth) -- (goth2);

\node[annot, align=center] at (-7.4, 4.4) {\textbf{Region Revenue Nodes}\\\textit{Americas, EMEA, APAC}};

% Stacked the Weak/Conditional label so it doesn't cross the line
\draw[dotedge, ->] (geo_hub.south) -- (-2.4, -2.8) node[midway, left=2pt, font=\scriptsize, align=right] {Weak/Conditional\\Link};


% -------------------------------------------------------------------------
% D. BOTTOM BRANCH: PRODUCT HIERARCHY [2, 3] & WRDS BUS FINANCIALS [9]
% -------------------------------------------------------------------------
\node[hubnode] (prod_hub) at (0, -1.8) {Product\\Hierarchy};

\node[segnode] (prod_sa) at (-2.0, -3.4) {$\text{Segment}_a$};
\node[segnode] (prod_sn) at (0.0, -3.8) {$\text{Segment}_n$};
\node[segnode] (prod_sb) at (2.0, -3.4) {$\text{Segment}_b$};

% Dashed oval grouping the segments for Financial Attributes
\node[draw=black!85, dashed, thick, rounded corners=12pt, inner sep=6pt, fit=(prod_sa) (prod_sn) (prod_sb)] (fin_box) {};

% MOVED WRDS Business Financials label much further right (to x=6.8) so the edge text has plenty of horizontal room
\node[annot, font=\footnotesize\bfseries, align=center] (fin_label) at (7.8, -3.6) {WRDS Business\\Financials [9]};
\draw[->, thick, black!85] (fin_box.east) -- node[midway, above=0.1pt, font=\scriptsize] {\texttt{node\_role = OPERATING}} (fin_label.west);

% Product/Service Leaf Nodes
\node[prodleaf] (leaf_pa) at (-2.0, -5.2) {Atomic Nodes [3]\\$a, d, e$};
\node[prodleaf] (leaf_pn) at (0.9, -5.5) {Atomic Nodes [3]\\$a, b, n$};
\node[prodleaf] (leaf_pb) at (3.8, -5.2) {Atomic Nodes [3]\\$c, f, g$};

\draw[edgestyle] (comp8) -- (prod_hub);
\draw[edgestyle] (prod_hub) -- (prod_sa) node[midway, right=60pt, font=\tiny] {Segment Nodes [2]};
\draw[edgestyle] (prod_hub) -- (prod_sn);
\draw[edgestyle] (prod_hub) -- (prod_sb);
\draw[edgestyle] (prod_sa) -- (leaf_pa);
\draw[edgestyle] (prod_sn) -- (leaf_pn);
\draw[edgestyle] (prod_sb) -- (leaf_pb);

% =========================================================================
% 2. MACRO-LEVEL TOPOLOGY (BOTTOM: FLEXIBLE NETWORK C_1 ... C_108)
% =========================================================================

% Well-spaced Spiral of Macro Nodes ordered by Market Cap
\node[macronode, minimum size=0.42cm, font=\tiny] (c1) at (-0.6, -9.2) {$C_1$};
\node[macronode, minimum size=0.50cm, font=\tiny] (c2) at (-1.5, -9.1) {$C_2$};
\node[macronode, minimum size=0.58cm, font=\scriptsize] (c3) at (-1.7, -10.1) {$C_3$};
\node[macronode, minimum size=0.66cm, font=\scriptsize] (c4) at (-0.6, -10.7) {$C_4$};
\node[macronode, minimum size=0.75cm, font=\scriptsize] (c5) at (0.5, -10.0) {$C_5$};
\node[macronode, minimum size=0.86cm, font=\footnotesize] (c6) at (0.4, -8.6) {$C_6$};
\node[macronode, minimum size=0.98cm, font=\footnotesize] (c7) at (-1.2, -7.8) {$C_7$};
\node[macronode, minimum size=1.18cm, fill=gray!28, double, double distance=1.3pt, thick] (c8) at (-3.1, -8.8) {$C_8$};
\node[macronode, minimum size=1.28cm] (c9) at (-2.3, -11.4) {$C_9$};
\node[macronode, minimum size=1.40cm] (c10) at (0.4, -11.7) {$C_{10}$};
\node[macronode, minimum size=1.52cm] (c11) at (2.5, -10.1) {$C_{11}$};
\node[macronode, minimum size=1.80cm] (c108) at (1.5, -7.4) {$C_{108}$};

% Subtle dotted spiral guide showing increasing Market Cap
\draw[black!40, semithick, dotted] (c1) -- (c2) -- (c3) -- (c4) -- (c5) -- (c6) -- (c7) -- (c8) -- (c9) -- (c10) -- (c11);

% Illustrative inter-firm connections formed from micro-graphs / attributes
\draw[black!85, semithick] (c8) -- (c7);
\draw[black!85, semithick] (c8) -- (c3);
\draw[black!85, semithick] (c8) -- (c9);
\draw[black!85, semithick] (c7) -- (c108);
\draw[black!85, semithick] (c6) -- (c108);
\draw[black!85, semithick] (c5) -- (c10);
\draw[black!85, semithick] (c4) -- (c9);

% --- Clean Unobstructed Zoom Arrow ---
\draw[->, very thick, black!90] (c8.north west) .. controls (-4.4, -5.0) and (-3.0, -1.2) .. (comp8.south west);

% --- Simple Side Annotations ---
\node[annot] (cap_note) at (5.8, -10.8) {Node size represents market cap;\\firms also carry 10-K filing date};
\draw[->, semithick, black!80] (3.4, -10.8) -- (cap_note.west);

\node[annot] (conn_note) at (6.0, -8.2) {\textbf{TEMPORAL\_SUCCESSOR} edges\\track longitudinal segment rewiring};
\draw[black!75, semithick] (3.5, -7.3) -- (3.5, -9.1);

\end{tikzpicture}%
}
\caption{Schematic of the two-layer Hierarchical Network of Networks $\Hnet$. \textbf{Bottom (Macro-Level Topology $\Gmacro$):} The static cohort of $N=108$ QQQ firms ($C_1, \dots, C_{108}$) arranged in increasing order of annual market capitalization (indicated by node size) and annotated with firm-level attributes including Form 10-K filing dates. Macro-level connections form flexibly from the firms' attributes and underlying micro-graphs. \textbf{Top (Micro-Level Topology $\Gmicro$):} Zoom-in on representative firm \texttt{Company 8} ($C_8$), showing the 5 core subgraphs spanning the 10-table schema: (1) Product Hierarchy [2, 3], (2) Market Competition [4, 5], (3) Dual-Scope Risk Taxonomy [6, 7, 8], (4) WRDS Business Financials [9], and (5) WRDS Geographic Financials [10], which are connected via solid structural edges, dashed weak links, and M:N associative bridges.}
\label{fig:micro_macro_schema}
\end{figure}

\subsection{Micro-Level Topology: Intra-Firm Sub-Networks}
For each corporation $C_i$ ($i \in \{1, \dots, 108\}$) in year $t$, the intra-firm topology is modeled as a rooted, multi-branch hierarchical graph:
\begin{equation}
    \Gmicro = \left( V_i^{(t)}, \; E_i^{(t)}, \; \mathbf{A}_i^{(t)} \right)
\end{equation}
centered on the company root vertex $r_i$, which stores firm-level metadata including the company's Central Index Key (CIK), annual market capitalization $M(C_i^{(t)})$, and exact \textbf{Form 10-K filing date} $d_{\text{file}}(C_i^{(t)})$. As depicted in the top panel of Fig.~\ref{fig:micro_macro_schema}, the vertex set $V_i^{(t)}$ is partitioned into four hierarchical tiers radiating outward from $r_i$:
\begin{equation}
    V_i^{(t)} = \{r_i\} \cup V_{i,\text{hub}} \cup V_{i,\text{inter}}^{(t)} \cup V_{i,\text{leaf}}^{(t)}
\end{equation}

\begin{enumerate}
    \item \textbf{Primary Category Hubs ($V_{i,\text{hub}}$):} Four first-tier functional hubs connect directly to the company root $r_i$:
    \begin{equation*}
        V_{i,\text{hub}} = \left\{ v_{\text{prod}}, \; v_{\text{comp}}, \; v_{\text{risk}}, \; v_{\text{geo}} \right\}
    \end{equation*}
    representing \texttt{Product/Services}, \texttt{Competitor}, \texttt{Risk}, and \texttt{Geographic \& Other Revenue}, respectively.
    
    \item \textbf{Intermediate Segment and Partition Nodes ($V_{i,\text{inter}}^{(t)}$):} The second structural tier organizes disclosures by operating segment and firm-wide scope:
    \begin{itemize}
        \item Under the \texttt{Product/Services} hub ($v_{\text{prod}}$), edges branch directly into the firm's active operating segments $S_i^{(t)} = \{\text{Segment}_a, \text{Segment}_b, \dots, \text{Segment}_n\}$. Each segment node $s \in S_i^{(t)}$ carries a quantitative \textbf{Financial Attribute} $w_{\text{rev}}(s) \in \mathbf{A}_i^{(t)}$ corresponding to its reported segment revenue from Items~7 and~8.
        \item Under both \texttt{Competitor} ($v_{\text{comp}}$) and \texttt{Risk} ($v_{\text{risk}}$), the hierarchy first bifurcates into a \texttt{Segment} partition node---which subsequently branches into the segment set $S_i^{(t)}$---and a general partition node (\texttt{Others} for unsegmented enterprise-wide competitors; \texttt{Other / External} for macroeconomic, regulatory, or systemic risks).
        \item Under \texttt{Geographic \& Other Revenue} ($v_{\text{geo}}$), the branch bifurcates into \texttt{Geographic} and \texttt{Other} revenue partition nodes.
    \end{itemize}

    \item \textbf{Entity Leaf Nodes ($V_{i,\text{leaf}}^{(t)}$):} The terminal layer contains the granular operational and geographic entities extracted in Tranche~3:
    \begin{itemize}
        \item \textit{Product/Service Leaves:} Granular offerings (e.g., $\text{Prod/Services}_{a,d,e}$) connected to one or more parent segments in $S_i^{(t)}$, permitting cross-segment product sharing.
        \item \textit{Competitor and Risk Leaves:} Specific named competitors and granular risk factors (depicted as shaded circular leaf nodes in Fig.~\ref{fig:micro_macro_schema}) attached either to their corresponding operating segment $\text{Segment}_k$ or to the \texttt{Others} / \texttt{Other/External} branch.
        \item \textit{Geographic Region Leaves:} Region-level nodes ($\text{Region}_a, \text{Region}_b, \dots, \text{Region}_n$, representing individual countries or broader continental regions depending on how the filer reports geographic disaggregation) and an aggregated $\text{Rest}$ node under the \texttt{Geographic} branch, annotated with revenue attributes from Items~7 and~8.
    \end{itemize}
\end{enumerate}

\subsection{Macro-Level Topology: Attribute- and Micro-Graph-Driven Connectivity}
At the macroscopic scale, the inter-firm network in year $t$ is defined over the fixed node set $\Vmacro = \{C_1, C_2, \dots, C_{108}\}$:
\begin{equation}
    \Gmacro = \left( \Vmacro, \; \Emacro, \; \mathbf{A}_{\text{macro}}^{(t)} \right)
\end{equation}
Each macro-node $C_i \in \Vmacro$ is endowed with macro-level attributes $\mathbf{a}_i^{(t)} \in \mathbf{A}_{\text{macro}}^{(t)}$, notably its annual market capitalization $M(C_i^{(t)})$ and its \textbf{Form 10-K filing date} $d_{\text{file}}(C_i^{(t)})$ (capturing fiscal calendar alignment and reporting timing). As illustrated in the bottom panel of Fig.~\ref{fig:micro_macro_schema}, macro-nodes in each snapshot $t$ are visually arranged in a spiral of increasing market capitalization:
\begin{equation}
    M\!\left(C_1^{(t)}\right) \le M\!\left(C_2^{(t)}\right) \le \dots \le M\!\left(C_{108}^{(t)}\right)
\end{equation}
where node size is proportional to $M(C_i^{(t)})$.

Crucially, rather than restricting $\Gmacro$ to a single fixed connectivity rule, our framework allows the macro-edge set $\Emacro$ to emerge flexibly from any chosen combination of firm-level attributes $\mathbf{A}_{\text{macro}}^{(t)}$ and the underlying micro-networks $\{\Gmicro\}_{i=1}^{108}$:
\begin{equation}
    e_{ij}^{(t)} = f\!\left(\mathbf{a}_i^{(t)}, \, \mathbf{a}_j^{(t)}, \, \Gmicro, \, G_{j,\text{micro}}^{(t)}\right)
\end{equation}
where $f(\cdot)$ represents the selected relational mapping---such as shared risk exposures, direct competitor citations, product/sector overlap, geographic revenue co-exposure, or reporting-period alignment via 10-K filing dates. Because pairwise relationships are induced directly from empirical firm attributes and micro-graph structures across all $\binom{108}{2} = 5,778$ pairs, the macro-topology naturally adapts to the analytical dimension of interest while allowing firms with no structural overlap under a given criterion (e.g., node $C_{11}$ in Fig.~\ref{fig:micro_macro_schema}) to remain unconnected.